\documentclass[11pt]{article}
\usepackage[T1]{fontenc}
\usepackage{textcomp}
\usepackage[margin=0.75in]{geometry}
\usepackage{amsmath,amssymb,amsthm}
\usepackage{array,booktabs}
\usepackage{hyperref}
\setlength{\parindent}{0pt}
\newtheorem{theorem}{Theorem}
\theoremstyle{definition}
\newtheorem{definition}{Definition}
\title{Flow Proof Artifact: prop-27-triangle-half-parallelogram}
\author{Generated by \texttt{flow doc proof}}
\date{Flow Proof Book}
\begin{document}
\maketitle
A triangle on the same base and between the same parallels equals half the parallelogram.\\[0.5em]
\textbf{Source.} euclid — Elements, Book I, Proposition 27\\[0.5em]
\bigskip
\noindent{\large \textbf{Derived fact 1.} \textit{A triangle on the same base and between the same parallels equals half the parallelogram} \hfill the Euclidean plane \cdot Euclid Book I \cdot Proposition 27: a triangle is half a parallelogram on same base}\par
\smallskip\noindent\textit{Coordinate: the Euclidean plane \cdot Euclid Book I \cdot Proposition 27: a triangle is half a parallelogram on same base}\par
\smallskip\noindent\textit{Source: euclid — Elements, Book I, Proposition 27}\par
\smallskip\noindent\textit{Built on: proposition 25: parallelograms on the same base and parallels are equal, for Book I of the Elements in on the Euclidean plane, proposition 4: side-angle-side congruence, for Book I of the Elements in on the Euclidean plane}\par
\medskip
\noindent\textbf{Goal.} A triangle on the same base and between the same parallels equals half the parallelogram\par
\noindent$\displaystyle \triangle ABC = \text{half parallelogram ABCD}$\par
\medskip
\noindent
\begin{tabular}{@{} >{\raggedright\arraybackslash}p{0.06\textwidth} >{\raggedright\arraybackslash}p{0.47\textwidth} >{\raggedleft\arraybackslash}p{0.41\textwidth} @{}}
\textbf{\#} & \textbf{Proof} & \textbf{Mathematics} \\
\midrule
\textbf{1.} & We prove that proposition 27: a triangle is half a parallelogram on same base for Book I of the Elements in the Euclidean plane. & \\[0.45em]
\textbf{2.} & Let diagonal\_AC = segment\_from\_A\_to\_C. & \\[0.45em]
\textbf{3.} & We invoke the derived fact: triangle\_ABC\_and\_parallelogram\_ABCD\_on\_same\_base\_and\_parallels. & \\[0.45em]
\textbf{4.} & We invoke the derived fact governing Book I of the Elements in the Euclidean plane: proposition 25: parallelograms on the same base and parallels are equal, for Book I of the Elements in on the Euclidean plane. & \\[0.45em]
\textbf{5.} & We invoke the axiom governing Book I of the Elements in the Euclidean plane: proposition 4: side-angle-side congruence, for Book I of the Elements in on the Euclidean plane. & \\[0.45em]
\textbf{6.} & From step 2, step 3, step 4, and step 5, by the chain of results established in Book I (step 2, step 3, step 4, and step 5), we obtain triangle ABC equals half parallelogram ABCD. Hence proven. & $\displaystyle \triangle ABC = \text{half parallelogram ABCD}$ \\[0.45em]
\bottomrule
\end{tabular}
\label{thm:Geometry-Euclid Book I-proposition_27_a_triangle_is_half_a_parallelogram_on_same_base}
\par\medskip\hrule
\end{document}
